% §2 Background and Threat Model --- Target 2--2.5pp

\section{Background and Threat Model}
\label{sec:background}

\subsection{Covenants and the Four Introspection Primitives}
\label{sec:bg:covenants}

A \emph{covenant} is a spending condition that restricts not only who
can spend a UTXO but how the resulting transaction is
structured~\cite{MES16}. Implementing one requires \emph{transaction
introspection}: the ability of a script to examine fields of the
spending transaction. Every covenant proposal we study reduces to a
choice of \emph{introspection primitive}: the mechanism by which the
script observes transaction data. The four choices correspond to axis
\axintro{} in Section~\ref{sec:bg:axes} and determine each proposal's
downstream attack surface.

\paragraph{Digest equality (\opctv{}, BIP-119~\cite{BIP119}).} The
script carries a $32$-byte template hash. At spend time, the consensus
engine computes SHA-$256$ over the canonical concatenation
$\mathtt{nVersion}\,\|\,\mathtt{nLockTime}\,\|\,\mathtt{scriptSigsHash}\,\|\,\mathtt{inputCount}\,\|\,\mathtt{sequenceHash}\,\|\,\mathtt{outputCount}\,\|\,\mathtt{outputsHash}\,\|\,\mathtt{inputIndex}$
and checks equality against the committed hash. The primitive is
whole-template: any field divergence (including \texttt{outputsHash},
which itself commits to the \texttt{(amount,\,scriptPubKey)} pair of
every output) produces digest mismatch and rejection.

\paragraph{Merkle-membership with amount semantics (\opccv{},
BIP-443~\cite{Ing23}).} \opccv{} takes a mode byte selecting which
per-output check applies: $\{-1,\,0,\,1,\,2\}$ for self-input check,
output script check, output amount deduction, and combined checks
respectively. Recursion is supported via taptree preservation, which
enables partial unvault with automatic revault. Any other mode byte
triggers \texttt{OP\_SUCCESSx} and silently disables the covenant
check~(axis \axopsurf{}, threat TM8).

\paragraph{Dedicated opcodes (\opvault{}, BIP-345~\cite{BIP345}).}
\texttt{OP\_VAULT} (trigger) and \opvaultrecover{} (recovery) are
purpose-built opcodes with vault-specific semantics fixed at the
consensus layer. Amount preservation, automatic revault, and recovery
authorisation via a \texttt{recoveryauth} Schnorr key are enforced by
the opcode rather than by script composition. BIP-345 was withdrawn in
early 2025 in favour of BIP-443~\cite{OB25}; we retain it as the
purpose-built design point and a direct economic comparator for
\opccv{}.

\paragraph{Sighash-preimage reconstruction (\catcsfs{},
BIP-347+348~\cite{Poe21,BIP347,BIP348}).} \texttt{OP\_CAT} concatenates stack
items; \opcsfs{} verifies a Schnorr signature against a
stack-supplied message. The vault assembles the BIP-341 sighash
preimage from witness fields, then verifies the same signature under
the same public key against both the assembled preimage (via
\opcsfs{}) and the consensus-computed sighash (via
\texttt{OP\_CHECKSIG}). Soundness reduces to Schnorr's EUF-CMA under
the fixed-public-key model~\cite{NRS20}: a valid signature over two
distinct messages under the same key would imply a solution to the
discrete logarithm problem. Consequently the assembled preimage must
equal the consensus preimage, which binds the output set through
\texttt{hashOutputs}. The vault uses
\sighash{SINGLE}$|$\texttt{ACP} so that external fee-paying inputs
may be attached without breaking the binding. The recovery leaf,
however, is a bare \texttt{cold\_pk~OP\_CHECKSIG} with no
introspection, leaving the recovery destination
output-unconstrained~(axis \axrecbind{}, threat TM11).

\subsection{The Four Bitcoin Script Extension Proposals}
\label{sec:bg:designs}

Each proposal is a specific composition of the primitives above with
key-gating and fee-inclusion decisions. The dimensions that matter for
the analysis in Section~\ref{sec:impossibility} are enumerated as a
closed set of axes in Section~\ref{sec:bg:axes} below, and each
proposal's position on every axis is the canonical
Table~\ref{tab:design_space}.

\paragraph{\opctv{}.} The vault uses a two-step deterministic
transaction tree: a bare \opctv{} script guards the deposit, and the
unvault UTXO offers a hot withdrawal after a CSV delay or a
\opctv{}-committed cold sweep. An anchor output sized at $3{,}300$
sats supports CPFP fee bumping because the committed template cannot
anticipate future fee rates (axis \axfee{} =
\texttt{out-of-band-anchor}).

\paragraph{\opccv{}.} Per-output granularity enables partial unvault
with revault; the recovery leaf is \emph{keyless} (axis \axauth{} =
\texttt{keyless}), which creates a permissionless-griefing surface
(TM2) and is also the sole axis value that admits class \classgrief{}
as susceptible.

\paragraph{\opvault{}.} A three-key model separates the trigger key,
the \texttt{recoveryauth} key, and the recovery scriptPubKey.
\opvault{}'s value-preservation rule forces a separate fee-wallet
input on every trigger and recovery transaction, adding $80$--$90$\vB{}
of overhead but retaining dynamic fee inclusion (axis \axfee{} =
\texttt{dynamic}).

\paragraph{\catcsfs{}.} The dual-verification pattern (Section~\ref{sec:bg:covenants})
gives the strongest hot-key binding of the four designs (axis
\axwdbind{} = \texttt{signature-dual-bind}) but pairs it with the
weakest recovery binding (axis \axrecbind{} =
\texttt{plain-signature}); the cold-key
compromise surface (TM11) follows mechanically.

\subsection{The Design-Space Axes}
\label{sec:bg:axes}

The four proposals differ along seven axes. Three carry over from the
formal decomposition~\cite{SHMB20,Har24} and appear with Greek-letter
symbols ($f$, $a$, $g$, $b$) in proofs; the remaining four are
introduced in this work and are used as descriptive labels. Every
attack class in Section~\ref{sec:impossibility} derives from a
specific axis-value; every design's vulnerability profile is a lookup
against Table~\ref{tab:design_space}.

\begin{description}[leftmargin=1.5em,itemsep=1pt,parsep=0pt,topsep=2pt]
  \item[\axauthf{} --- Recovery authorisation.] Who may initiate the
    recovery path? Values: \texttt{keyless}, \texttt{recoveryauth-key},
    \texttt{hot-key}, \texttt{cold-key}. Controls class \classgrief{}
    (permissionless griefing).
  \item[\axrevaultf{} --- Revault support.] Does the covenant allow
    partial withdrawal? Values: \texttt{yes}, \texttt{no}. Controls
    class \classscale{} (per-UTXO recovery-cost scaling).
  \item[\axfeef{} --- Fee-inclusion model.] Where does the spending
    transaction's fee come from? Values: \texttt{dynamic} (fee in the
    transaction), \texttt{static-precommit} (fee baked at creation),
    \texttt{out-of-band-anchor} (CPFP via a separate anchor key).
    Controls class \classfee{} (fee-channel pinning).
  \item[\axwdbind{} --- Withdrawal-destination binding.] How is the
    withdrawal destination constrained? Values:
    \texttt{template-hash-digest}, \texttt{per-output-amount+script},
    \texttt{dedicated-opcode}, \texttt{signature-dual-bind}. Controls
    class \classhot{} (hot-key theft via destination substitution).
  \item[\axrecbind{} --- Recovery-destination binding.] Same question
    for the recovery path. Values: same as \axwdbind{}, plus
    \texttt{plain-signature} (no script-level constraint). Controls
    class \classcold{} (cold-key theft via unconstrained recovery).
    Kept separate from \axwdbind{} because \catcsfs{} occupies
    different positions on the two sub-axes --- strongest on
    \axwdbind{}, weakest on \axrecbind{}.
  \item[\axopsurf{} --- Opcode surface.] Is the covenant expressed by
    a single semantic opcode or a multi-mode polymorphic one with
    \texttt{OP\_SUCCESSx} fallthrough on undefined modes? Values:
    \texttt{single-mode}, \texttt{multi-mode-with-OP\_SUCCESS}.
    Controls class \classmode{} (mode/parameter bypass).
  \item[\axintro{} --- Introspection primitive.] How does the script
    read transaction data? Values correspond to the four primitives of
    Section~\ref{sec:bg:covenants}:
    \texttt{digest-equality},
    \texttt{merkle-membership-and-amount},
    \texttt{dedicated-opcode}, and
    \texttt{sighash-preimage-reconstruction}. Descriptive; governs
    which of the other axes are simultaneously realisable.
\end{description}

For the formal decomposition we retain $(f,\,a,\,g,\,b)$ as the axes
of the design-space lattice used in Propositions~1 and~2: $f =
\axfee$, $a = \axrevault$, $g = \axauth$, and $b$ is the joint
$(\axwdbind,\,\axrecbind)$ sublattice. Axes \axopsurf{} and \axintro{}
are additional dimensions orthogonal to the proofs; they are where the
developer-footgun (class \classmode{}) and the design-implementability
questions live.

\begin{table}[t]
\centering
\caption{The four Bitcoin covenant proposals, positioned on
the seven axes of Section~\ref{sec:bg:axes}. Every attack-class
susceptibility claim in Section~\ref{sec:impossibility} is a lookup
against this table. Simplicity (Appendix~D) occupies the strict-best
position on \axauth{}, \axrevault{}, \axwdbind{}, \axrecbind{}, and
\axintro{} simultaneously; no Bitcoin proposal does.}
\label{tab:design_space}
\footnotesize
\setlength{\tabcolsep}{2.2pt}
\renewcommand{\arraystretch}{1.12}
\begin{tabular}{@{}lllll@{}}
\toprule
Axis                                    & \opctv{}                       & \opccv{}                               & \opvault{}                    & \catcsfs{}                      \\
\midrule
\axauth{}  (recovery auth)              & hot-key                        & keyless                                & recoveryauth-key              & cold-key                        \\
\axrevault{} (revault)                  & no                             & yes                                    & yes                           & no                              \\
\axfee{}   (fee model)                  & out-of-band-anchor             & dynamic                                & dynamic                       & dynamic                         \\
\axwdbind{} (withdrawal binding)        & template-hash-digest           & per-output-amount+script               & dedicated-opcode              & signature-dual-bind             \\
\axrecbind{} (recovery binding)         & template-hash-digest           & per-output-amount+script               & dedicated-opcode              & plain-signature                 \\
\axopsurf{}  (opcode surface)           & single-mode                    & multi-mode-with-OP\_SUCCESS            & single-mode                   & single-mode                     \\
\axintro{}   (introspection primitive)  & digest-equality                & merkle-membership-and-amount           & dedicated-opcode              & sighash-preimage-reconstruction \\
\bottomrule
\end{tabular}
\renewcommand{\arraystretch}{1.0}
\end{table}

Two structural dependencies between axes constrain the realisable
combinations and are used implicitly throughout
Section~\ref{sec:impossibility}: (i) \axfee{} =
\texttt{out-of-band-anchor} implies \axrevault{} = \texttt{no}, because
the committed template cannot anticipate arbitrary future splits; and
(ii) \axintro{} = \texttt{sighash-preimage-reconstruction} is the only
known primitive that admits \axwdbind{} =
\texttt{signature-dual-bind}, because the binding reduces to Schnorr
message-equality under a single key.

\subsection{Vault Lifecycle, Threat Taxonomy, and Threat Model}
\label{sec:bg:lifecycle}

All four designs implement the Swambo et al.\ lifecycle~\cite{SHMB20}:
a deposit UTXO is triggered into an unvaulting UTXO subject to a
relative timelock (CSV); after the delay, the hot key completes the
withdrawal; at any point, the cold key (or, for \opccv{}, any party)
may invoke emergency recovery to a pre-committed address. Security
depends on an online watchtower detecting unauthorised triggers within
the CSV window.

\paragraph{Threat model (Kerckhoffs).} We adopt the standard
cryptographic threat model. The adversary knows the vault's design,
the covenant type, the script structure, and all public parameters
--- including the recovery address where applicable. Security rests on
keys and on the script's semantics, not on the secrecy of
construction. Every attack-feasibility claim in this work is evaluated
under this assumption. Concretely, we do not rely on taproot leaf
hiding, address non-disclosure, or closed-source construction as
security mechanisms.

\medskip We extend Swambo et al.\ with design-specific attack classes.
The 11-element threat taxonomy below serves as citation anchors for
the per-class analysis of Section~\ref{sec:impossibility} (classes
\classfee{}--\classmode{} subsume groups of TMs; see
Table~\ref{tab:design_space} and Appendix~B for the full class$\times$TM
map):

{\footnotesize
\begin{description}[leftmargin=1.3em,itemsep=0pt,parsep=0pt,topsep=2pt]
  \item[TM1 (fee-channel pinning).] Hot-key attacker saturates Bitcoin
    Core's $25$-descendant mempool cap with dust-rate transactions
    spending the \opctv{} unvault's anchor, blocking the watchtower's
    CPFP recovery. Class \classfee{}.
  \item[TM2 (recovery griefing).] An unauthorised party---anyone,
    under keyless recovery; the cold or auth-key holder, under
    key-gated---repeatedly invokes recovery, cycling the owner
    through re-deposit and re-trigger. Class \classgrief{}.
  \item[TM3 (trigger-key theft).] Hot-key attacker initiates an
    unauthorised unvault and races the watchtower during the CSV
    delay. Class \classhot{}.
  \item[TM4 (per-UTXO recovery-cost scaling).] Partial unvault with
    revault grows the vault's UTXO set over time; per-UTXO recovery
    cost scales linearly in the count and in fee rate. Under
    congestion, some UTXOs fall below the economic abandonment
    threshold. An attacker with the trigger key can accelerate this
    distribution; ordinary operation can produce it without one.
    Previously framed as ``watchtower exhaustion''; the reframe
    separates the structural cost tail from the multi-round attacker
    narrative (see Section~\ref{sec:empirical:batching} and
    Appendix~B).
    Class \classscale{}.
  \item[TM5 (trigger-xpub theft).] Capture of \opvault{}'s trigger
    xpub amplifies TM3 to every vault derived from that xpub.
  \item[TM6 (dual-key compromise).] Simultaneous theft of both the
    hot/trigger key and the cold/\texttt{recoveryauth} key defeats
    the split-key safety assumption.
  \item[TM7 (address reuse).] A second deposit to the same \opctv{}
    vault address produces a UTXO whose amount does not match the
    committed template hash, rendering it permanently unspendable. An
    operational-correctness failure rather than an adversarial attack
    (no adversary is required); included here as an axis-\axwdbind{}
    consequence observable under normal wallet use.
  \item[TM8 (\opccv{} mode bypass).] An undefined mode byte in the
    \opccv{} witness triggers \texttt{OP\_SUCCESSx}, silently
    disabling the covenant check. Specified behaviour under
    BIP-443's forward-compatibility design (axis \axopsurf{} =
    \texttt{multi-mode-with-OP\_SUCCESS}); the contribution is the
    systematic mode sweep and class-level framing, not the discovery.
    Class \classmode{}.
  \item[TM9 (hot-key redirection).] A \catcsfs{} hot-key attempt to
    redirect the withdrawal output is rejected by the dual
    \opcsfs{}+\texttt{OP\_CHECKSIG} binding; the outcome degrades to
    griefing only. Immunity derivation for class \classhot{}.
  \item[TM10 (witness tampering).] Modification of \catcsfs{}
    sighash-preimage witness fields is detected by the
    dual-verification check (stack-assembled preimage must agree with
    the consensus-computed sighash). Cryptographic verification,
    not an attack class.
  \item[TM11 (cold-key theft).] Capture of the cold/recovery key
    grants unrestricted fund redirection if the recovery path is
    output-unconstrained (\catcsfs{}); forced-recovery-only if
    output-bound (\opctv{}, \opvault{}). Class \classcold{}.
\end{description}
}

Three threats dominate the analysis: TM1 (\opctv{} anchor CPFP), TM2
(\opccv{} keyless recovery), and TM4 (per-UTXO recovery-cost scaling
on \opccv{}/\opvault{}). The per-design susceptibility matrix appears
in Section~\ref{sec:imposs:design_space}.

\providecommand{\todo}[2][]{\textbf{[TODO: #2]}}
